\section{Verification of the Concurrent Code}
\label{sec:verification}

The engine holds no lock, so its correctness under concurrency rests on memory-ordering arguments.
A stress test samples only a small fraction of the possible interleavings: a determinism gate of
150 runs explores almost none of the schedules its thread count admits. The concurrent code is
therefore checked by two model checkers in addition to the test suite.

\paragraph{Model checking under weak memory.}
GenMC~\cite{kokologiannakis2021genmc} enumerates every execution of a bounded concurrent program
under the RC11 memory model~\cite{lahav2017repairing}: for the given thread count, operation count
and inputs, it explores every interleaving and every choice of which write each read observes. A
clean result therefore covers all executions within the bound. Three harnesses are weaker. Two
bound the number of context switches, which the checker supports only under sequential
consistency, so they cover the sequentially consistent executions. The third samples executions
because its shape is too large to enumerate, and the same property is checked exhaustively at a
smaller shape by another harness.

Each harness is a small program that includes the engine's headers and calls its functions, so it
checks the engine's code rather than a separate model, and it stops compiling when the code it
checks changes shape. Two harnesses check protocols inside functions the checker cannot load, and
contain verbatim copies of that code running over the engine's structures.

The harnesses check the following properties, each of which would fail without an error:
\begin{itemize}
\item \emph{Concurrent map:} get-or-create has exactly one winner and returns the winner's value to
every caller, including when a resize runs concurrently with the inserts that caused it and across
two successive resizes; the ``inserted'' flag comes from the publishing exchange, so repeated
offers of one value are not each reported as insertions.
\item \emph{Key set:} two threads inserting one key while a third forces a resize produce exactly
one reported insertion; an enumeration during a resize reports each key once; and a membership
query for a key inserted before the query began answers present, including while the key is
being copied to the new table. This last property requires walking the chain of tables from the
oldest; the harness contains a twenty-execution counterexample for newest-first order.
\item \emph{Work-stealing deque:} the owner and a thief never take the same item, checked where the
two ends address the same slot.
\item \emph{Match rendezvous:} a match is claimed exactly once and never dropped on a collision,
and every pair of an instance and a match recorded for its class is replayed.
\item \emph{Per-instance application list:} of two applications linked into one list, exactly one
sees the other, so each branchial pair is emitted once (Section~\ref{sec:quotient}).
\item \emph{Arena:} no two workers receive the same arena index, and the shared allocation path
never hands out memory a worker's private cursor owns.
\item \emph{Scheduler:} a worker that waits is always woken. Removing either of the two
sequentially consistent fences in the submit and wait protocol makes the checker report a
non-terminating spin loop.
\end{itemize}

Six harnesses check the GPU kernel's protocols under the scoped extension of the memory model: the
ring buffer's exactly-once hand-off, termination detection without early exit, the ring and the
termination detector together, the hash table electing one inserter, the replay rendezvous meeting
every pair, and the causal dynamic program's producer--transition rendezvous.

Harnesses that link the whole engine construct an engine and add a rule; at the stated bound the
checker reaches the end of each (442 complete executions), so the composed construction is covered
exhaustively.

To run the engine through GenMC, the checker was modified; the patch against release 0.17.0 is
shipped with the engine. It has four correctness fixes (\texttt{operator new[]} treated as an
allocation; typed promotion of \texttt{memcpy} and \texttt{memset}; the failure ordering for a
failed compare-exchange's read; and no liveness report while a thread is blocked by an assume,
upstream pull request 58), a second pass of memory-intrinsic promotion after scalar replacement,
and instrumentation: a progress report, and an attribution of the state-space estimate to the
source lines whose reads and writes produce it.

\paragraph{TLA\textsuperscript{+} models.}
Four protocols are also specified in TLA+ and checked with its model checker~\cite{lamport2002specifying}.
Every protocol step in progress is held in a set from which any enabled step may fire, so the
checked interleavings cover every worker count; the bound is on the number of states instead. Each
specification is a transcription of the code and names the code it transcribes. Each has a broken
configuration that the checker must reject, which shows that the specification can detect the
failure it is about.

\emph{Match forwarding} is the property whose failure costs the most: a lost match removes the
whole subtree below it and leaves the run internally consistent. The specification is checked in
four configurations (Table~\ref{tab:tla}), which differ in whether the process that wins a match's
claim both stores and propagates it, and in whether submission is batched. The shipped
configuration is complete at the bound. Without ownership and without batching the property fails
on a chain $s_0 \to s_1 \to s_2$: a match is found at the root after $s_2$'s pull has run; $s_1$'s
pull wins the claim and stores the match without propagating it; the root's push finds the claim
taken and does not recurse, so $s_2$ never receives the match. Batching alone hides this failure at
the bound, so ownership is required independently of batching.

\begin{table}[htbp]
\centering
\caption{\textbf{Forwarding completeness under TLA+}, exhaustive at the bound of four states, three
matches and three edges, with a discovery mid-chain that enables a grandchild. \emph{Ownership}:
the claim winner both stores and propagates the match; \emph{batched}: a state's matches are
submitted together after its matching completes.}
\label{tab:tla}
\resizebox{\ifdim\width>\linewidth\linewidth\else\width\fi}{!}{%
\begin{tabular}{lllcr}
\toprule
Configuration & Ownership & Batched & Verdict & Distinct states \\
\midrule
\textbf{shipped} & yes & yes & complete       & 117{,}005 \\
calibration & no  & yes & complete       & 85{,}777 \\
calibration & yes & no  & complete       & 479{,}005 \\
calibration & no  & no  & violated (as required) & trace depth 13 \\
\bottomrule
\end{tabular}}
\end{table}

\emph{Depth relaxation} (Section~\ref{sec:depth}): at quiescence the set of expanded states is
exactly the set whose shortest-path depth is below the budget, including when a steered step lowers
the budget for one pass. The broken configurations derive a child's depth from its parent's depth
at claim time, and claim states against a bound different from the one the match task is admitted
under; both violate the property.

\emph{Termination} (Section~\ref{sec:depth}): the job system never reports completion while work
remains, provided a job submits all of its work before it is counted as complete. The broken
configuration, in which a job submits work after it has been counted, violates the property.

\emph{Segmented array:} the array that stores states, events, edges and application lists grows by
segments, and if a segment exists then every segment below it exists. The shipped protocol is
clean at $2{,}284$ distinct states; allocating only the requested segment violates the property.
Three appenders are needed to expose the failure: with two, each claims once, and the index that
would reveal the gap cannot be reached before the lower segment is published. GenMC cannot
construct this structure, so it is checked in TLA+.

\paragraph{Sanitizers.}
The test suite runs under ThreadSanitizer, AddressSanitizer and UndefinedBehaviorSanitizer, and the
device suites under the CUDA memory and initialization checkers; all five report nothing and fail
the build on a finding. The initialization checker catches host reads of device memory that no
kernel wrote, which return whatever the allocator last left there, so a test on such a value passes
or fails depending on allocation history.
